\documentclass[11pt,a4paper]{article}

% =========================================================================
%  RAPPORT SAE CYBER — Conception réseau sécurisé GNS3
%  Theme : Réseau d'entreprise — palette indigo / cyan
% =========================================================================

\usepackage[utf8]{inputenc}
\usepackage[T1]{fontenc}
\usepackage[french]{babel}
\usepackage{lmodern}
\usepackage[a4paper,margin=2.2cm,top=2.6cm,bottom=2.4cm]{geometry}
\usepackage{xcolor}
\usepackage{fontawesome5}
\usepackage{tikz}
\usetikzlibrary{shapes.geometric,positioning,calc,arrows.meta}
\usepackage{titlesec}
\usepackage[most]{tcolorbox}
\usepackage{listings}
\usepackage{fancyhdr}
\usepackage{booktabs}
\usepackage{tabularx}
\usepackage{array}
\usepackage{colortbl}
\usepackage{enumitem}
\usepackage{microtype}
\usepackage[bookmarks=true,colorlinks=true,linkcolor=accent,urlcolor=accent,citecolor=accent]{hyperref}

% --- Palette indigo cyan -------------------------------------------------
\definecolor{primary}{HTML}{1E1B4B}
\definecolor{accent}{HTML}{4F46E5}
\definecolor{accent2}{HTML}{06B6D4}
\definecolor{accentdark}{HTML}{312E81}
\definecolor{light}{HTML}{EEF2FF}
\definecolor{midgray}{HTML}{6B7280}
\definecolor{codebg}{HTML}{1A1B3D}
\definecolor{codefg}{HTML}{F0F2FF}
\definecolor{codekw}{HTML}{93C5FD}
\definecolor{codestr}{HTML}{A5F3FC}
\definecolor{codecmt}{HTML}{7C84BD}

\titleformat{\section}{\sffamily\Large\bfseries\color{primary}}{\thesection}{0.8em}{}[\vspace{-0.4ex}{\color{accent}\rule{\linewidth}{1.2pt}}]
\titleformat{\subsection}{\sffamily\large\bfseries\color{accent}}{\thesubsection}{0.6em}{}
\titleformat{\subsubsection}{\sffamily\normalsize\bfseries\color{primary}}{\thesubsubsection}{0.5em}{}
\titlespacing*{\section}{0pt}{1.6em}{0.8em}

\pagestyle{fancy}
\fancyhf{}
\renewcommand{\headrulewidth}{0pt}
\fancyhead[L]{\footnotesize\sffamily\color{midgray}\faProjectDiagram~SAE Cyber — Réseau sécurisé GNS3}
\fancyhead[R]{\footnotesize\sffamily\color{midgray}Honorine \& Grondin}
\fancyfoot[C]{\footnotesize\sffamily\color{midgray}\thepage}

\lstdefinestyle{cisco}{
  backgroundcolor=\color{codebg},
  basicstyle=\ttfamily\footnotesize\color{codefg},
  keywordstyle=\color{codekw}\bfseries,
  stringstyle=\color{codestr},
  commentstyle=\color{codecmt}\itshape,
  showstringspaces=false, breaklines=true,
  frame=none, framesep=8pt, framerule=0pt,
  xleftmargin=12pt, xrightmargin=12pt,
  columns=fullflexible, inputencoding=utf8, extendedchars=true,
  morekeywords={vlan,interface,switchport,access,trunk,allowed,encapsulation,mode,name,ip,address,router,ospf,network,area,standby,priority,preempt,access-list,permit,deny,extended,mpls,bgp,neighbor,remote-as,update-source,activate,address-family,vpnv4,send-community,exit-address-family,route-target,export,import,Loopback0,GigabitEthernet,no,shutdown,hostname,Vlan,spanning-tree,extend,system-id},
  literate={é}{{\'e}}1 {è}{{\`e}}1 {ê}{{\^e}}1 {à}{{\`a}}1 {ô}{{\^o}}1 {î}{{\^\i}}1 {ç}{{\c{c}}}1 {ù}{{\`u}}1 {â}{{\^a}}1 {É}{{\'E}}1 {È}{{\`E}}1 {Ê}{{\^E}}1 {À}{{\`A}}1 {Ô}{{\^O}}1 {Ç}{{\c{C}}}1 {Ù}{{\`U}}1 {Â}{{\^A}}1
}
\lstset{style=cisco}

\newtcblisting{ciscobox}[1][]{
  listing only,
  enhanced jigsaw, breakable,
  colback=codebg, colframe=accent,
  arc=2pt, boxrule=0pt, leftrule=3pt,
  left=8pt, right=8pt, top=4pt, bottom=4pt,
  listing options={style=cisco}, #1
}

\newtcolorbox{infobox}[1][]{
  enhanced, colback=light, colframe=accent,
  arc=2pt, boxrule=0pt, leftrule=3pt,
  left=10pt, right=10pt, top=6pt, bottom=6pt,
  fontupper=\small, #1
}
\newtcolorbox{techbox}[1][]{
  enhanced, colback=accent2!10, colframe=accent2,
  arc=2pt, boxrule=0pt, leftrule=3pt,
  left=10pt, right=10pt, top=6pt, bottom=6pt,
  fontupper=\small, #1
}

\setlist[itemize]{leftmargin=*,itemsep=2pt,topsep=4pt}
\setlist[enumerate]{leftmargin=*,itemsep=2pt,topsep=4pt}

\begin{document}
\pagenumbering{gobble}

\begin{tikzpicture}[remember picture, overlay]
  \fill[primary] (current page.north west) rectangle ([yshift=-9cm]current page.north east);
  \fill[accent] ([yshift=-9cm]current page.north west) rectangle ([yshift=-9.25cm]current page.north east);
  % Logo : reseau interconnecté (network nodes)
  \node[anchor=center] at ([yshift=-4.6cm]current page.north) {%
    \begin{tikzpicture}[scale=0.95]
      % Site 1
      \fill[white] (-2.4,1.2) circle (0.45);
      \node[scale=1.0,accent] at (-2.4,1.2) {\faServer};
      % Site 2
      \fill[white] (2.4,1.2) circle (0.45);
      \node[scale=1.0,accent] at (2.4,1.2) {\faServer};
      % Site 3
      \fill[white] (0,-1.6) circle (0.45);
      \node[scale=1.0,accent] at (0,-1.6) {\faServer};
      % central PE
      \fill[accent2] (0,0.2) circle (0.65);
      \node[scale=1.4,white] at (0,0.2) {\faNetworkWired};
      % liens
      \draw[white,line width=2pt] (-2.4,1.2) -- (-0.5,0.45);
      \draw[white,line width=2pt] (2.4,1.2) -- (0.5,0.45);
      \draw[white,line width=2pt] (0,-1.2) -- (0,-0.4);
      % petits noeuds
      \foreach \x/\y in {-3.1/0.4, -1.7/2.1, 3.1/0.4, 1.7/2.1, -0.8/-1.95, 0.8/-1.95} {
        \fill[accent!60] (\x,\y) circle (0.15);
      }
      \foreach \a/\b/\c/\d in {-3.1/0.4/-2.6/1.05, -1.7/2.1/-2.25/1.55, 3.1/0.4/2.6/1.05, 1.7/2.1/2.25/1.55, -0.8/-1.95/-0.35/-1.78, 0.8/-1.95/0.35/-1.78} {
        \draw[accent!60,line width=0.8pt] (\a,\b) -- (\c,\d);
      }
    \end{tikzpicture}};
  \node[anchor=north, white, font=\sffamily\Huge\bfseries] at ([yshift=-6.7cm]current page.north) {Réseau Sécurisé};
  \node[anchor=north, accent2, font=\sffamily\Large] at ([yshift=-7.65cm]current page.north) {Conception multi-sites sous GNS3};
  \node[anchor=north, white!75!primary, font=\sffamily\small] at ([yshift=-8.4cm]current page.north) {VLAN \textbullet{} HSRP \textbullet{} MPLS VPN \textbullet{} ACL};

  \node[anchor=south west, primary, font=\sffamily\small] at ([xshift=2.2cm,yshift=3cm]current page.south west) {\textbf{\textcolor{accent}{Auteurs}}};
  \node[anchor=south west, primary, font=\sffamily\small] at ([xshift=2.2cm,yshift=2.55cm]current page.south west) {HONORINE Kylian};
  \node[anchor=south west, primary, font=\sffamily\small] at ([xshift=2.2cm,yshift=2.25cm]current page.south west) {GRONDIN Angélique};
  \node[anchor=south east, primary, font=\sffamily\small] at ([xshift=-2.2cm,yshift=3cm]current page.south east) {\textbf{\textcolor{accent}{Encadrant}}};
  \node[anchor=south east, primary, font=\sffamily\small] at ([xshift=-2.2cm,yshift=2.5cm]current page.south east) {SITAL DAHONE Aloïs};
  \fill[accent] ([yshift=1.3cm]current page.south west) rectangle ([yshift=1.4cm]current page.south east);
  \node[anchor=south, midgray, font=\sffamily\footnotesize] at ([yshift=0.6cm]current page.south) {IUT Réseaux \& Télécommunications — Filière Cybersécurité};
\end{tikzpicture}
\newpage

\pagenumbering{arabic}
\setcounter{page}{1}

{\sffamily
\hypersetup{linkcolor=primary}
\renewcommand{\contentsname}{\textcolor{accent}{\faList~} Sommaire}
\tableofcontents
}
\vspace{0.4cm}

\section{Introduction}

\begin{infobox}[title={\faBullseye~Mission SAE 3.Cyber 03}]
Concevoir l'infrastructure réseau sécurisée d'une entreprise multi-sites, en interconnectant trois sites distants via un \textbf{VPN MPLS} tout en garantissant sécurité, redondance et fiabilité.
\end{infobox}

\noindent Les briques techniques mises en \oe{}uvre :
\begin{itemize}
  \item \textbf{VLANs} pour isoler logiquement les services (RH, Vente, AdminRéseau, etc.)
  \item \textbf{Routage inter-VLAN} sur switches niveau 3 (SD1, SD2)
  \item \textbf{HSRP} pour la redondance des passerelles
  \item \textbf{PVST+} pour l'optimisation du Spanning Tree par VLAN
  \item \textbf{ACL} pour le contrôle d'accès granulaire
  \item \textbf{MPLS VPN} avec MP-BGP pour l'interconnexion sécurisée des sites
\end{itemize}

Validation sous \textbf{Packet Tracer} et \textbf{GNS3}, en suivant les recommandations de l'ANSSI.

% =========================================================================
\section{Configuration des switchs L2}

\subsection{Switch SA1 — accès aux VLANs métiers}

Les VLANs sont déclarés par service (RH, Vente, Admin, Serveur fichier, Client Wi-Fi) et chaque port d'accès est rattaché au VLAN qui correspond :

\begin{ciscobox}
vlan 10
 name RH
vlan 20
 name Vente
vlan 30
 name AdminReseau
vlan 40
 name serverfichier
vlan 50
 name clientwifi

interface Ethernet0/1
 switchport access vlan 10
 switchport mode access

interface Ethernet0/2
 switchport access vlan 20
 switchport mode access

interface Ethernet0/0
 switchport trunk allowed vlan 10,20,30,40,50,60,99
 switchport trunk encapsulation dot1q
 switchport mode trunk
\end{ciscobox}

\subsection{Switches SA2 et SA3 — VLAN unique}

Ces switches ne portent qu'un seul VLAN (60 — \og other \fg{}) destiné aux VPCS. Leur configuration est volontairement minimaliste :

\begin{ciscobox}
vlan 60
 name other

interface Ethernet0/1
 switchport access vlan 60

interface Ethernet0/0
 switchport trunk allowed vlan 60
 switchport trunk encapsulation dot1q
 switchport mode trunk
\end{ciscobox}

% =========================================================================
\section{Routage inter-VLAN sur les switches L3}

\subsection{SD1 — routage des VLANs 10 à 50 + 99}

Chaque VLAN métier reçoit une interface SVI dotée de sa passerelle. Le VLAN 99 est dédié au management :

\begin{ciscobox}
interface Vlan10
 ip address 192.168.10.1 255.255.255.0
interface Vlan20
 ip address 192.168.20.1 255.255.255.0
interface Vlan30
 ip address 192.168.30.1 255.255.255.0
interface Vlan40
 ip address 192.168.40.1 255.255.255.0
interface Vlan50
 ip address 192.168.50.1 255.255.255.0
interface Vlan99
 ip address 192.168.99.1 255.255.255.0

interface gi0/0
 switchport access vlan 99
interface gi0/1
 switchport trunk allowed vlan 10,20,30,40,50
\end{ciscobox}

\subsection{SD2 — VLAN 60 + management}

\begin{ciscobox}
interface Vlan60
 ip address 192.168.60.1 255.255.255.0
interface Vlan99
 ip address 192.168.99.2 255.255.255.0

interface gi0/1
 switchport trunk allowed vlan 60
interface gi0/0
 switchport access vlan 99
\end{ciscobox}

% =========================================================================
\section{Haute disponibilité — HSRP et PVST+}

\subsection{HSRP — passerelle virtuelle redondée}

Le protocole \textbf{HSRP} (Hot Standby Router Protocol) permet à deux routeurs/switches L3 de partager une IP virtuelle. En cas de panne du primaire, le secondaire prend le relais sans interruption pour les clients.

\textbf{SD1 (actif, priorité 110) :}
\begin{ciscobox}
interface Vlan10
 ip address 192.168.10.2 255.255.255.0
 standby 10 ip 192.168.10.1
 standby 10 priority 110
 standby 10 preempt

interface Vlan20
 ip address 192.168.20.2 255.255.255.0
 standby 20 ip 192.168.20.1
 standby 20 priority 110
 standby 20 preempt
! ... idem Vlan30, 40, 50, 60, 99
\end{ciscobox}

\textbf{SD2 (standby, priorité par défaut 100) :}
\begin{ciscobox}
interface Vlan10
 ip address 192.168.10.3 255.255.255.0
 standby 10 ip 192.168.10.1
 standby 10 preempt
! ... idem pour les autres VLANs
\end{ciscobox}

\begin{techbox}[title={\faSync~Comportement observé}]
À l'arrêt de SD1, le trafic bascule automatiquement vers SD2 sans rupture côté client. Les IPs physiques sont en \texttt{.2} (SD1) et \texttt{.3} (SD2), l'IP virtuelle exposée aux clients en \texttt{.1}.
\end{techbox}

\subsection{PVST+ — Spanning Tree par VLAN}

Le mode \texttt{pvst} (Per-VLAN Spanning Tree) permet d'avoir une instance STP par VLAN, ce qui ouvre la porte à un équilibrage de charge fin entre liens :

\begin{ciscobox}
spanning-tree mode pvst
spanning-tree extend system-id
spanning-tree vlan 10 priority 4096    ! SD1 sera root pour VLAN 10
\end{ciscobox}

Sur SD2, on choisit un autre VLAN (ex. 20) comme root pour exploiter le second lien physique \textemdash{} le trafic de chaque VLAN emprunte un chemin différent en régime nominal.

% =========================================================================
\section{Contrôle d'accès — ACL}

\subsection{ACL RH (sur SD1) — protection du serveur RH}

Le serveur RH (\texttt{192.168.10.10}) n'est accessible en HTTPS qu'à partir de Vente (entièrement) et d'AdminRéseau (un seul poste, \texttt{192.168.30.30}) :

\begin{ciscobox}
ip access-list extended ACL_RH
 permit tcp 192.168.20.0 0.0.0.255 host 192.168.10.10 eq 443
 permit tcp host 192.168.30.30      host 192.168.10.10 eq 443
 deny   tcp any                     host 192.168.10.10 eq 443
 permit ip  any                     host 192.168.10.10

interface vlan 10
 ip access-group ACL_RH in
\end{ciscobox}

\subsection{ACL serveur de fichiers (sur SD1)}

Le serveur FTP (\texttt{192.168.40.40}) accepte FTP (port 21) et data (port 20) depuis Vente et le poste admin uniquement :

\begin{ciscobox}
ip access-list extended ACL_FTP
 permit tcp 192.168.20.0 0.0.0.255 host 192.168.40.40 eq 21
 permit tcp 192.168.20.0 0.0.0.255 host 192.168.40.40 eq 20
 permit tcp host 192.168.30.30      host 192.168.40.40 eq 21
 permit tcp host 192.168.30.30      host 192.168.40.40 eq 20
 deny   tcp any                     host 192.168.40.40 eq 21
 deny   tcp any                     host 192.168.40.40 eq 20
 permit ip  any                     host 192.168.40.40
\end{ciscobox}

\subsection{ACL 104 — isolation du poste admin}

Le poste \texttt{192.168.30.30} ne doit être joignable que depuis le VLAN d'administration (99) :

\begin{ciscobox}
access-list 104 permit ip 192.168.99.0 0.0.0.255 host 192.168.30.30
access-list 104 deny   ip 192.168.10.0 0.0.0.255 host 192.168.30.30
access-list 104 deny   ip 192.168.20.0 0.0.0.255 host 192.168.30.30
access-list 104 deny   ip 192.168.30.0 0.0.0.255 host 192.168.30.30
access-list 104 deny   ip 192.168.40.0 0.0.0.255 host 192.168.30.30
access-list 104 deny   ip 192.168.50.0 0.0.0.255 host 192.168.30.30
access-list 104 permit ip any any
\end{ciscobox}

% =========================================================================
\section{Pré-VPN MPLS — pont entre SD et PE}

Les routeurs PE étant en mode router, ils ne peuvent pas faire de \texttt{switchport mode access}. La solution retenue : déclarer un VLAN \texttt{1xx} entre chaque SD et son PE, et faire passer OSPF dessus pour échanger les réseaux.

\textbf{Côté SD (exemple SD1, VLAN 100) :}
\begin{ciscobox}
interface vlan 100
 ip address 192.168.100.1 255.255.255.252

interface GigabitEthernet1/0
 switchport access vlan 100
 switchport trunk allowed vlan 10,20,30,40,50,60,99,100
 switchport trunk encapsulation dot1q
 switchport mode access

router ospf 1
 network 192.168.10.0  0.0.0.255 area 0
 network 192.168.20.0  0.0.0.255 area 0
 ! ... idem 30, 40, 50
 network 192.168.100.0 0.0.0.3   area 0
\end{ciscobox}

\textbf{Côté PE :}
\begin{ciscobox}
interface vlan 100
 ip address 192.168.100.2 255.255.255.252

router ospf 1
 network 192.168.100.0 0.0.0.3 area 0
\end{ciscobox}

% =========================================================================
\section{Cœur du projet — VPN MPLS multi-sites}

\subsection{Principe}

Trois sites, chacun avec deux switches L3 (SDx). Un routeur \textbf{PE central} (PE4) agrège les flux MPLS via MP-BGP VPNv4. Chaque couple site-SD reçoit une \textbf{VRF} dédiée pour assurer l'étanchéité entre clients.

\subsection{PE1 — site 1}

\begin{ciscobox}
hostname PE1
mpls ldp router-id Loopback0 force

ip vrf site1sd1
 rd 100:1
 route-target export 100:1
 route-target import 100:1
ip vrf site1sd2
 rd 100:2
 route-target export 100:2
 route-target import 100:2

interface Loopback0
 ip address 1.1.1.1 255.255.255.255

interface GigabitEthernet0/0
 ip vrf forwarding site1sd1
 ip address 192.168.99.1 255.255.255.0
 mpls ip
 no shutdown
interface GigabitEthernet0/1
 ip vrf forwarding site1sd2
 ip address 192.168.99.2 255.255.255.0
 mpls ip
 no shutdown
interface GigabitEthernet0/2
 ip address 10.1.1.1 255.255.255.252
 mpls ip
 no shutdown

router bgp 100
 bgp router-id 1.1.1.1
 bgp log-neighbor-changes
 neighbor 10.1.1.4 remote-as 100
 neighbor 10.1.1.4 update-source Loopback0
 address-family vpnv4
  neighbor 10.1.1.4 activate
  neighbor 10.1.1.4 send-community both
 exit-address-family
\end{ciscobox}

\subsection{PE2 et PE3 — sites 2 et 3}

Configuration identique à PE1, en ajustant :
\begin{itemize}
  \item Loopback : \texttt{2.2.2.2} (PE2) / \texttt{3.3.3.3} (PE3)
  \item VRFs : \texttt{site2sd3}, \texttt{site2sd4} pour PE2 ; \texttt{site3sd5}, \texttt{site3sd6} pour PE3
  \item Route Distinguishers et Route Targets : \texttt{100:3}, \texttt{100:4}, \texttt{100:5}, \texttt{100:6}
  \item Liens vers PE central : \texttt{10.1.2.1/30} et \texttt{10.1.3.1/30}
\end{itemize}

\subsection{PE central (PE4) — agrégateur}

Le routeur central termine toutes les sessions MP-BGP et porte l'ensemble des VRFs. Il s'occupe aussi du routage OSPF sur les liens vers les PEs :

\begin{ciscobox}
hostname PE
mpls ldp router-id Loopback0 force

ip vrf site1sd1
 rd 100:1
 route-target export 100:1
 route-target import 100:1
! ... idem pour site1sd2, site2sd3, site2sd4, site3sd5, site3sd6

interface Loopback0
 ip address 4.4.4.4 255.255.255.255

interface GigabitEthernet0/0
 ip address 10.1.1.4 255.255.255.252
 mpls ip
interface GigabitEthernet0/1
 ip address 10.1.2.4 255.255.255.252
 mpls ip
interface GigabitEthernet0/2
 ip address 10.1.3.4 255.255.255.252
 mpls ip

router ospf 1
 router-id 4.4.4.4
 network 10.1.1.0 0.0.0.3 area 0
 network 10.1.2.0 0.0.0.3 area 0
 network 10.1.3.0 0.0.0.3 area 0

router bgp 100
 bgp router-id 4.4.4.4
 neighbor 10.1.1.1 remote-as 100
 neighbor 10.1.1.1 update-source Loopback0
 neighbor 10.1.2.1 remote-as 100
 neighbor 10.1.2.1 update-source Loopback0
 neighbor 10.1.3.1 remote-as 100
 neighbor 10.1.3.1 update-source Loopback0
 address-family vpnv4
  neighbor 10.1.1.1 activate
  neighbor 10.1.1.1 send-community both
  neighbor 10.1.2.1 activate
  neighbor 10.1.2.1 send-community both
  neighbor 10.1.3.1 activate
  neighbor 10.1.3.1 send-community both
 exit-address-family
\end{ciscobox}

% =========================================================================
\section{Architecture finale}

\renewcommand{\arraystretch}{1.35}
\begin{tabularx}{\linewidth}{@{}>{\bfseries}p{2.6cm} X@{}}
\toprule
\rowcolor{light} Couche & Rôle \\ \midrule
SA1, SA2, SA3   & Switches d'accès — port \textit{access} vers les hôtes \\
SD1, SD2 et co. & Switches L3 — routage inter-VLAN, HSRP, PVST+ \\
PE1, PE2, PE3   & Routeurs P/E des sites distants — MPLS LDP + MP-BGP VRF \\
PE (central)    & Concentrateur MPLS — reflector RR pour les VRFs \\
\bottomrule
\end{tabularx}

\section{Annexe — Debian + FTP sous GNS3}

\subsection{Préparation de l'image Debian}
\begin{enumerate}
  \item Télécharger l'appliance Debian sur \href{https://gns3.com/marketplace/appliances/debian-2}{gns3.com/marketplace} (fichier \texttt{.gns3a}).
  \item Récupérer l'image QEMU \texttt{debian-12.6.qcow2} sur SourceForge.
  \item Dans GNS3 : \texttt{New template > Import an appliance template file}, sélectionner le \texttt{.gns3a}, puis importer le \texttt{.qcow2}.
  \item Allouer 2-4\,Go de RAM et 2-3 c\oe{}urs CPU dans \texttt{Préférences > QEMU > Debian > General}.
\end{enumerate}

\subsection{Adressage statique}

\begin{ciscobox}[listing options={style=cisco,language=bash}]
# /etc/network/interfaces
auto ens5
iface ens5 inet dhcp

auto ens4
iface ens4 inet static
    address 192.168.40.40
    netmask 255.255.255.0
    gateway 192.168.40.1
    dns-nameservers 192.168.1.1
\end{ciscobox}

\subsection{Serveur FTP (vsftpd anonyme)}

\begin{ciscobox}[listing options={style=cisco,language=bash}]
apt update && apt install -y vsftpd ftp
\end{ciscobox}

Extrait clé de \texttt{/etc/vsftpd.conf} :
\begin{ciscobox}[listing options={style=cisco,language=bash}]
listen=YES
listen_ipv6=NO
anonymous_enable=YES
local_enable=NO
write_enable=YES
anon_upload_enable=YES
anon_mkdir_write_enable=YES
anon_root=/var/ftp
allow_writeable_chroot=YES
pasv_enable=YES
pasv_min_port=10000
pasv_max_port=10100
\end{ciscobox}

\begin{ciscobox}[listing options={style=cisco,language=bash}]
mkdir -p /var/ftp && chmod 755 -R /var/
chmod a-w /var/ftp
systemctl restart vsftpd
ftp localhost   # user: anonymous, pass: (vide)
\end{ciscobox}

\subsection{Route par défaut côté client}

\begin{ciscobox}[listing options={style=cisco,language=bash}]
ip route add default via 192.168.50.1 dev ens4
\end{ciscobox}

% =========================================================================
\section{Conclusion}

Le projet matérialise une infrastructure d'entreprise complète : segmentation par VLAN, routage L3, redondance HSRP, équilibrage PVST+, sécurité par ACL, et \textbf{interconnexion multi-sites par MPLS VPN}. L'agrégation via un PE central simplifie l'extension du réseau : ajouter un site revient à provisionner une VRF et une session BGP supplémentaires.

\medskip
\textbf{Avantages de l'architecture :}
\begin{itemize}
  \item \textbf{Simplicité de gestion} — un point central pour tout le routage inter-sites
  \item \textbf{Scalabilité} — chaque nouveau site = une VRF + une session BGP
  \item \textbf{Sécurité et isolation} — séparation stricte des clients par VRF
\end{itemize}

\medskip
\textbf{Pistes d'amélioration :} ajout de mécanismes de résilience BGP multi-path, VRRP pour les sites internes, durcissement des bornes selon les recommandations ANSSI.

\vspace{0.5cm}
\begin{center}
{\sffamily\itshape\color{midgray}\large
\og La VRF n'est pas qu'une table de routage : c'est une frontière. \fg{}}
\end{center}

\end{document}
